\documentclass[12pt]{article}
\usepackage[letterpaper,margin=0.85in]{geometry}
\usepackage{fontspec}
\setmainfont{Times New Roman}
\usepackage{enumitem}
\usepackage{array}
\usepackage{booktabs}
\usepackage{microtype}
\usepackage{titlesec}
\usepackage{fancyhdr}

\setlength{\parindent}{0pt}
\setlength{\parskip}{7pt}
\setlist{nosep,leftmargin=1.6em}
\titleformat{\section}[block]{\bfseries\large}{}{0pt}{}
\titlespacing*{\section}{0pt}{13pt}{4pt}
\pagestyle{fancy}
\fancyhf{}
\renewcommand{\headrulewidth}{0pt}
\fancyfoot[L]{\small CareBot three-bin design}
\fancyfoot[R]{\small \thepage}

\begin{document}

\begin{center}
    {\LARGE\bfseries CareBot Three-Bin Patient Delivery Robot}\par
    \vspace{5pt}
    {\normalsize Short design document for the first robot}
\end{center}

\section*{Basic idea}

CareBot is basically a robotic car that carries ten Red Cross patient cubes. We load the cubes before the run. The robot then follows a fixed path and drops each group at the correct place.

The robot has three bins. The two small bins hold two cubes each. The larger middle bin holds six cubes. Each bin has a flap underneath it and its own servo motor. When a servo turns, the flap opens and the cubes fall out.

\begin{center}
\renewcommand{\arraystretch}{1.35}
\begin{tabular}{|>{\centering\arraybackslash}p{0.22\textwidth}|>{\centering\arraybackslash}p{0.36\textwidth}|>{\centering\arraybackslash}p{0.22\textwidth}|}
\hline
Small bin A & Large middle bin & Small bin B \\
2 cubes & 6 cubes & 2 cubes \\
Servo flap 1 & Servo flap 2 & Servo flap 3 \\
\hline
\end{tabular}
\end{center}

The proposed outer size is 230 mm wide, 240 mm long, and 280 mm high. These are maximum values. The final bin sizes depend on the real cube size. We will therefore test the layout with cardboard first.

\section*{Movement sequence}

\begin{enumerate}
    \item The robot moves straight along the main path.
    \item At the first marker, it stops and opens small bin A. Two cubes drop.
    \item It turns left onto the 45-degree branch.
    \item It follows the next path using the line sensor and wheel encoders.
    \item At the middle marker, it opens the large bin. All six cubes drop.
    \item It continues to the final marker and opens small bin B. The last two cubes drop.
    \item The robot closes the flap and stops.
\end{enumerate}

The phrase ``ten steps forward'' sounds simple, but it will not be very accurate. Battery level and wheel slip can change the distance. Encoder counts or field markers will actually give us a more repeatable stop.

\newpage
\section*{Electronics}

An Arduino Uno R3 controls the drive motors, line sensor, wheel encoders, and three flap servos. One DRV8833 board controls the two drive motors. Two SG90 servos can run the small flaps. An MG90S is safer for the large flap because it carries six cubes.

The new design has only three servos. Hence, the Uno can control them directly. A PCA9685 is optional and can be added later if we add more servos. The servos still need a separate regulated 5 V supply. They should not draw their power through the Uno.

\section*{Camera and AI}

The ESP32-CAM identifies the colour of a cylindrical patient. It sends a short result, such as RED, GREEN, YELLOW, or UNKNOWN, to the Uno. The Uno then chooses the saved route for that colour.

The saved ``coordinates'' are actually line markers and encoder distances. The camera does not know the robot's exact position on the whole field. If the camera is unsure, the robot stops and checks again.

The three bins complete the preloaded cube task. The camera can identify loose cylindrical patients, but it cannot pick them up. If the robot must carry those patients too, we will add a simple front scoop or gripper.

\par\medskip
{\large\bfseries Build order}\par\smallskip

First, we will test one flap with the correct load. Then we will build all three bins and check that every cube falls cleanly. After that, we will tune the path and stopping points. The camera comes last, once the basic robot works properly.

\end{document}
